Navigating complex command-line interface (CLI) syntax remains a challenge for developers, yet resolving this through cloud-based AI assistants introduces severe privacy risks in secure environments. This thesis addresses this trade-off by designing and implementing a local-first, CPU-efficient CLI command option retrieval system that operates entirely offline. Built in Rust using Hexagonal Architecture, the system combines a quantized bi-encoder embedding model (\texttt{embeddinggemma-300M-Q8\_0}) with a BM25 keyword search, backed by a Transition-Driven Command Assembly and Validation engine (STRS). To run efficiently on standard CPU hardware without GPU acceleration, the binary is compiled targeting musl libc with hardware-level SIMD optimizations. Technical evaluation ($N=609$ queries) demonstrates high resource efficiency, delivering an average execution latency of 563.21~ms and a stable memory footprint of 658.62~MB. In terms of retrieval, the model achieves 80.30\% accuracy in identifying the target command-line utility. When the correct utility is selected, the validation engine blocks 39.06\% of commands due to option conflicts; when an incorrect utility is selected, command generation is blocked in 53.33\% of cases due to option conflicts. While only 11.82\% of all queries result in a command generated for the correct utility with all required options matched, among commands generated for the correct utility, 24.16\% contain all required options, 51.01\% match a subset, and 24.83\% contain none. Globally, silently incorrect recommendations occur in 21.35\% of cases (9.20\% wrong tool generated, and 12.15\% correct tool with zero option matches). An exploratory usability study ($N=9$) indicates high user satisfaction, showing a clear learning effect as the median completion time fell from 603~seconds for the first task to 214~seconds for the third task, while demonstrating the system's potential to assist users with unfamiliar, proprietary tools. Ultimately, while retrieval accuracy is acceptable but not yet optimal, these findings demonstrate that local-first command retrieval is computationally feasible, offering a viable, privacy-preserving alternative to cloud-based assistants.
